% ============================================================
%  Mise en page conforme à la note pédagogique ESPRIT
%  (Times New Roman 12 pt, marges 2,5 cm, interligne 1,15,
%   texte justifié, alinéa 0,5 cm, pagination « x/N » en bas à droite)
% ============================================================
\usepackage{fontspec}
% Times New Roman chargée par ses fichiers (Windows) : le nom de famille seul donne parfois
% un faux gras en italique selon le cache de polices. Ailleurs (Linux, CI) : TeX Gyre Termes.
\IfFontExistsTF{timesbd.ttf}
  {\setmainfont{times.ttf}[BoldFont=timesbd.ttf,ItalicFont=timesi.ttf,BoldItalicFont=timesbi.ttf]}
  {\IfFontExistsTF{Times New Roman}{\setmainfont{Times New Roman}}
    {\setmainfont{texgyretermes}[Extension=.otf,UprightFont=*-regular,BoldFont=*-bold,
                                 ItalicFont=*-italic,BoldItalicFont=*-bolditalic]}}
\usepackage{polyglossia}
\setdefaultlanguage{french}

\usepackage[a4paper,margin=2.5cm,headheight=15pt]{geometry}
\usepackage{setspace}
\setstretch{1.15}
\setlength{\parindent}{0.5cm}
\setlength{\parskip}{4pt}

\usepackage{graphicx}
\graphicspath{{figures/}}
\usepackage{xcolor}
\usepackage{amssymb} % \checkmark dans les tableaux
\usepackage{float}
\usepackage{booktabs,tabularx,longtable,multirow,array}
\usepackage{enumitem}
\setlist{itemsep=2pt,topsep=4pt}
\usepackage{caption}
\captionsetup{font=small,labelfont=bf,labelsep=period}
\usepackage{lastpage}
\usepackage{fancyhdr}
\usepackage{titlesec}
\usepackage{tcolorbox}
\usepackage{pgfgantt}
\usepackage{pdfpages}
\usepackage[numbers,sort&compress]{natbib}
\usepackage{csquotes}
\usepackage[hidelinks]{hyperref}

% Code en ligne (\texttt) : coupure possible après « . », « / », « - », « _ » et « : »
% (noms de paquets, d'événements et de variables qui débordaient de la marge)
\ExplSyntaxOn
\tl_new:N \l__pfe_tt_tl
\cs_new_protected:Npn \__pfe_tt_coupures:n #1
  {
    \tl_set:Nn \l__pfe_tt_tl {#1}
    \tl_replace_all:Nnn \l__pfe_tt_tl { . } { . \allowbreak }
    \tl_replace_all:Nnn \l__pfe_tt_tl { / } { / \allowbreak }
    \tl_replace_all:Nnn \l__pfe_tt_tl { - } { - \allowbreak }
    \tl_replace_all:Nnn \l__pfe_tt_tl { \_ } { \_ \allowbreak }
    \__pfe_tt_deux_points:NV \l__pfe_tt_tl \c_colon_str
    \tl_use:N \l__pfe_tt_tl
  }
% « : » est une lettre en syntaxe expl3, d'où la constante \c_colon_str
\cs_new_protected:Npn \__pfe_tt_deux_points:Nn #1#2
  { \tl_replace_all:Nnn #1 {#2} { #2 \allowbreak } }
\cs_generate_variant:Nn \__pfe_tt_deux_points:Nn { NV }
\NewCommandCopy \pfetextttorigine \texttt
\RenewDocumentCommand \texttt { m } { \pfetextttorigine { \__pfe_tt_coupures:n {#1} } }
\ExplSyntaxOff
\setlength{\emergencystretch}{2em}

% Couleurs du document (changer ici la couleur principale, par ex. celle du logo de l'entreprise)
\definecolor{principale}{HTML}{1F6FEB}
\definecolor{accent}{HTML}{E8590C}

% Légendes : « Fig. 1 », « Tab. 1 » (note pédagogique)
\addto\captionsfrench{%
  \renewcommand{\figurename}{Fig.}%
  \renewcommand{\tablename}{Tab.}%
  \renewcommand{\listfigurename}{Liste des figures}%
  \renewcommand{\listtablename}{Liste des tableaux}%
  \renewcommand{\bibname}{Bibliographie et netographie}%
  \renewcommand{\contentsname}{Table des matières}}

% Titres en gras, taille décroissante, numérotés 1, 1.1, 1.1.1
\titleformat{\chapter}[display]{\bfseries\centering}{\Large Chapitre \thechapter}{8pt}{\LARGE}
\titlespacing*{\chapter}{0pt}{-20pt}{24pt}
\titleformat{\section}{\bfseries\Large}{\thesection}{0.2cm}{}
\titleformat{\subsection}{\bfseries\large}{\thesubsection}{0.2cm}{}
\titleformat{\subsubsection}{\bfseries\normalsize}{\thesubsubsection}{0.2cm}{}
\setcounter{secnumdepth}{3}
\setcounter{tocdepth}{2}

% Pagination : « x/N » en bas à droite
\fancypagestyle{corps}{\fancyhf{}\renewcommand{\headrulewidth}{0pt}%
  \fancyfoot[R]{\thepage/\pageref{LastPage}}%
  \fancyhead[R]{\small\textit{\leftmark}}}
\fancypagestyle{liminaire}{\fancyhf{}\renewcommand{\headrulewidth}{0pt}\fancyfoot[R]{\thepage}}
% Les débuts de chapitre (style plain) suivent le style courant
\newif\ifcorps \corpsfalse
\fancypagestyle{plain}{\fancyhf{}\renewcommand{\headrulewidth}{0pt}%
  \fancyfoot[R]{\ifcorps\thepage/\pageref{LastPage}\else\thepage\fi}}
\renewcommand{\chaptermark}[1]{\markboth{#1}{}}

% Repères de rédaction : mettre \afficherfalse pour la version finale
\newif\ifafficher \affichertrue
\newcommand{\todo}[1]{\ifafficher\textcolor{red}{\textbf{[À compléter : #1]}}\fi}

% Encadré pour les définitions (note pédagogique)
\newtcolorbox{definition}[1]{colback=principale!4,colframe=principale!60,fonttitle=\bfseries,title=Définition : #1}

% Figure à partir d'un fichier, ou cadre provisoire si le fichier n'existe pas encore
\newcommand{\insererfigure}[4][0.95]{% [largeur]{fichier}{légende}{label}
  \begin{figure}[H]\centering
    \IfFileExists{figures/#2}{\includegraphics[width=#1\linewidth]{#2}}%
      {\fbox{\parbox[c][5cm][c]{0.8\linewidth}{\centering\textcolor{red}{Figure à produire : \texttt{#2}}}}}
    \caption{#3}\label{#4}
  \end{figure}}

% Figure large sur une page en paysage (diagrammes de séquence, de classes, de déploiement)
\usepackage{pdflscape}
\newcommand{\insererfigurepaysage}[3]{% {fichier}{légende}{label}
  \begin{landscape}
  \begin{figure}[p]\centering
    \IfFileExists{figures/#1}{\includegraphics[width=\linewidth,height=0.82\textheight,keepaspectratio]{#1}}%
      {\fbox{\parbox[c][5cm][c]{0.8\linewidth}{\centering\textcolor{red}{Figure à produire : \texttt{#1}}}}}
    \caption{#2}\label{#3}
  \end{figure}
  \end{landscape}}
